\section{Projektziele}

\subsection{Muss-Ziele}

\begin{itemize}[itemsep=0.4pt]
    \item Die Lerninhalte in denn GitHub Repositories sollen automatisiert verarbeitet werden.
    \item Die JSON-Kommentarblöcke und die darin enthaltenen Metadaten sollen extrahiert und in eine einheitliche Struktur gebracht werden.
    \item Die Daten sollen als Knoten(Challenges) und Kanten(Abhängigkeiten) in Neo4j importiert werden.
    \item Die REST-API soll die Graphdaten, getrennt vom Frontend, über Endpunkte aufrufen.
    \item Das Frontend soll die empfangenen Graphdaten strukturiert und einheitlich bereitstellen.
\end{itemize}

\subsection{Soll-Ziele}

\begin{itemize}[itemsep=0.4pt]
    \item Challengefilterung des gewünschten Keywords.
    \item Detaillierte Challenge anzeige.
    \item Hover Darstellung der aktuellen Challenge mit Ansicht auf die Abhängende und der Nachfolgenden Aufgabe(Challenge).
    \item Bereitstellung einer Funktion zum aktualisieren der Daten, entweder manuel per Button oder mit einer Zeiplaneinstellung.
    \item Listenansicht Schalter, um die Challenges schneller einzusehen.
\end{itemize}

\subsection{Kann-Ziele}

\begin{itemize}[itemsep=0.4pt]
    \item Erweiterte Suchfunktionen der Challenges und Keywords.
    \item Statistiken der erledigten, gemerkten oder noch zu erledigenden Lerninhalten.
    \item Registry-, Benutzerrollen- und Authentifizierungsverwaltung per Keycloak.
\end{itemize}

\subsection{Projektabgrenzung}

\begin{itemize}[itemsep=0.4pt]
    \item Die eigentlichen Lerninhalte werden nicht erstellt oder überarbeitet.
    \item Eine mobile Lernanwendung ist nicht Bestandteil meines Projekts.
    \item Eine Skalierung für eine große Anzahl an parallelen Nutzern ist immoment nicht vorgesehen.
\end{itemize}